% !TEX TS-program = lualatex

%%%
% src::
%   url = https://github.com/pfradin/luadraw/discussions/80#discussioncomment-14370417
%%%

\documentclass[border = 3pt]{standalone}

%%%
% Des couleurs faciles d'emploi via le package \xcolor!
%%%
\usepackage[svgnames]{xcolor}

%%%
% La bibliothèque \luadraw allie une facilité d’utilisation
% à un rendu particulièrement soigné.
%%%
\usepackage{luadraw}

\begin{document}

\begin{luadraw*}{name = complementary-set}
local ld = luadraw

local cpx = ld.cpx

------
-- Nos ingrédients de base.
------
local i = cpx.I

local A = ld.set(0, 90)

local lab_A     = "$A$"
local lab_A_not = "$\\overline{A}$"
local lab_E     = "$\\Omega$"

local col_A     = "blue"
local col_A_not = "red"
local col_E     = "black"

------
-- ¨Def de la zone graphique.
------
local graphview = ld.graph:new{
  window = {-7.25, 7.25, -5, 5},
  size   = {10, 10},
  bbox   = false
}

------
-- Opacité des remplissages.
------
graphview:Fillopacity(0.2)

------
-- Comme la ¨doc de ¨luadraw nous informe que ''set(...)'' est
-- un chemin destiné aux commandes ''path'' et ''Dpath'', nous
-- allons construire nous-même l'univers et le complémentaire
-- de `A` via le langage de tracé proposé par les commandes
-- ''path'' et ''Dpath''. Voici des explications sur le code
-- utilisé ci-dessous.
--
--     1) ''Universe_and_A_not'' est une table contenant les
--     instructions de tracé.
--
--     1) ''{-6.5 + 4*i, ... , "cla"}'' sert au tracé de la
--     boîte arrondie resprésentant l'univers où ''"cla"'' est
--     pour "closed line arc", soit "une ligne fermé aux coins
--     arrondis"
--
--     1) ''A_bis'' est une copie des instructions définissant
--     ''A'' (pas besoin de refaire le travail abattu par la
--     commande ''set'' de ¨luadraw).
--
--     1) ''table.insert(A_bis, 2, "m")'' permet de déplacer le
--     tracé de `A` au point d'affixe complexe `2`. La commande
--     ''"m"'' est pour "move to", soit se "déplacer vers".
--
--     1) Lors du tracé, ''Dpath'', en lisant les instructions,
--     trace le rectangle, puis se déplace en `2` pour commencer
--     un nouveau tracé, à savoir celui de l'ensemble `A`.
--
--     1) Il ne reste plus qu'à utiliser la régle de remplissage
--     latex::''even odd rule'' de ¨tikz.
------
local Universe_and_A_not = {
  -6.5 + 4*i, 6.5 + 4*i, 6.5 - 4*i, -6.5 - 4*i, 0.25,
  "cla"
}

local A_bis = table.copy(A)
table.insert(A_bis, 2, "m")

ld.insert(Universe_and_A_not, A_bis)

graphview:Dpath(
  Universe_and_A_not,
     col_E .. ", fill = " .. col_A_not
  .. ", even odd rule"
)

------
-- Tracé de l'ensemble `A`.
------
graphview:Dpath(
  A,
  col_A .. ", fill = " .. col_A
)

------
-- Ajout des étiquettes.
------
graphview:Fillopacity(.75)

graphview:Dlabel(
  lab_A,
  0,
  {
    node_options = col_A,
    pos          = "center"
  },
  lab_A_not,
  4 + 2.7*i ,
  {
    node_options = col_A_not,
    pos          = "E"
  },
  lab_E,
  -6.5 + 4*i,
  {
    node_options = col_E,
    pos          = "NW"
  }
)

------
-- Montrons le résultat de notre oeuvre.
------
graphview:Show()
\end{luadraw*}

\end{document}
